\begin{python0}
from solutions import *; clear()
\end{python0}
\sectionthree{The boolean condition}

Here's one of our first few programs:
\begin{console}
std::cout << "entering the for-loop ...\n";
for (int i = 10; i > 0; i--)
{   
    std::cout << "entered the block ...\n";
    std::cout << "i: " << i;
    std::cout << "exiting the block ...\n";
}
std::cout << "done with the for-loop ...\n";
\end{console}

Look at this boolean expression:
\begin{console}[commandchars=\~\@\$]
std::cout << "entering the for-loop ...\n";
for (int i = 10; ~EMPHASIZE@i > 0$; i--)
{   
    std::cout << "entered the block ...\n";
    std::cout << "i: " << i;
    std::cout << "exiting the block ...\n";
}
std::cout << "done with the for-loop ...\n";
\end{console}
One of the most common misconceptions about the for-loop (and actually
applies to all loops) is that many first-time programmers think that if
the boolean expression
\begin{center}
\EMPHASIZE{i $>$ 0}
\end{center}
is false at \EMPHASIZE{any point} in the loop, you will break you out of
the for-loop immediately. That's \EMPHASIZE{NOT TRUE}. Look
at the picture on page 12 again. The check on whether to break out of
the loop occurs at a specific point. If I do this instead:
\begin{console}[commandchars=\~\@\$]
std::cout << "entering the for-loop ...\n";
for (int i = 10; ~EMPHASIZE@i > 0$; i--)
{   
    std::cout << "entered the block ...\n";
    std::cout << "i: " << i;
    ~EMPHASIZE@i = 0;$
    std::cout << "exiting the block ...\n";
}
std::cout << "done with the for-loop ...\n";
\end{console}


The program would still execute the last statement in the body of the
for-loop:
\begin{center}
\texttt{std::cout << "exiting the block ...\ n";}
\end{center}
because the boolean check comes after this print statement.
\begin{center}
\EMPHASIZE{MAKE SURE YOU REMEMBER THAT!}
\end{center}

